\section{Introduction}
\label{sec:intro}

A deployed language-model agent is a compound system. Around a frozen model sit a system prompt, tool descriptions, tool code, and increasingly multi-file skill packages \citep{agent_skills_docs}. Editing these files changes task success by tens of points without touching weights, and a growing family of optimizers searches for the edits: prompt and pipeline optimizers \citep{dspy,miprov2-dspy,gepa}, skill methods \citep{skillopt,evoskill,trace2skill,wikiskill}, context playbooks \citep{ace}, and harness search \citep{meta-harness,xu2026harnesslens}. Each paper reports a gain on its own benchmark, split, scorer, budget and optimizer model. The numbers cannot be lined up.

Our own history shows how this goes wrong. An earlier \capevolve{} draft reported a Gemma-4-31B gain of \histGemmaDeltaSaved{} points on SpreadsheetBench. When we recomputed it from the run artifacts, the same outputs scored \histGemmaDeltaRecalc{} points after formulas were recalculated, because the original scorer reads cached cell values and so rewards an agent for writing literals instead of formulas (Section~\ref{sec:problem}). The held-out split had been scored in earlier development runs. Optimizer sessions had read gold workbooks from a directory mounted next to the task inputs. None of this was visible in the headline table.

This paper is a measurement-first comparison. We fix the protocol before looking at any outcome and ask which optimizer finds the best artifact under it. The protocol has five parts: a functional scorer that is reported next to the legacy one; gold isolation, with a negative test for every agentic proposer; one promotion rule for the \capevolve{} arms, chosen from a budget-exact simulation (Section~\ref{sec:method:P}); a matched rollout checkpoint rather than a nominal budget (Section~\ref{sec:setup}); and final tasks the search never touched, scored once and analysed with paired, family-clustered statistics (Section~\ref{sec:setup:stats}). Figure~\ref{fig:pipeline} shows the pipeline.

Within that protocol we compare \capevolve{}, whose editable surface is a parameter of one loop with two search algorithms, against the official code of GEPA, SkillOpt, Meta-Harness, EvoSkill, Trace2Skill, ACE, MIPROv2 and HarnessLens, a community port of WikiSkill (no official code exists), and two controls, an unoptimized seed and a one-shot repair. The suites are SpreadsheetBench verified-400 (Suite A, Gemma-4-31B target), $\tau^2$-bench retail (Suite B) and airline (Suite C), with Claude Sonnet 5.5 as the optimizer for every matched arm. If a control or a simpler method wins, that is the main result.

Our contributions hold whichever arm wins.
\begin{enumerate}[leftmargin=*,nosep]
\item An audit of historical agent-optimization results: scorer compliance accounts for \histGemmaCompliancePts{} of \histGemmaDeltaSaved{} and \histGptCompliancePts{} of \histGptDeltaSaved{} reported points, and the replayed gate decisions of \replayRuns{} historical runs change in \replayFlipsUnmasked{} of \replayDecisions{} cases under a stricter rule (Section~\ref{sec:problem}).
\item A measurement protocol with its operating characteristics: the promotion rule has a false-promotion probability of \simPFalse{} in a null search of 12 candidates and a power of \simPPowerEight{} at a true gain of $+0.08$, which we state as underpowered by design (Section~\ref{sec:method:P}, Appendix~\ref{app:sim}).
\item A matched-checkpoint comparison of nine competing methods, two \capevolve{} algorithms and two controls over three suites, with full configurations, costs and infrastructure incidents reported (Section~\ref{sec:results}, Appendices~\ref{app:config} and~\ref{app:incidents}). \needs{Headline findings paragraph after the final evaluation.}
\end{enumerate}
